%%%%%%%%%%%%%%%%%%%%%%%%%%%%%%%%%%%%%%%%%%%%%%%%%%%%%%%%%%%%%%%%%%%%%%
% Lớp 11 | Đạo hàm | Đạo hàm
% Dạng 1D7?1-3: Ý nghĩa Hình học của đạo hàm
% Số câu: 27  (sinh tự động bởi tach_cau.py; từ nay thầy sửa trực tiếp file này)
%%%%%%%%%%%%%%%%%%%%%%%%%%%%%%%%%%%%%%%%%%%%%%%%%%%%%%%%%%%%%%%%%%%%%%

%%%%%%%%%%%%%%%% TRẮC NGHIỆM 4 PHƯƠNG ÁN %%%%%%%%%%%%%%%%

% nguồn: Ok_TracNghiemLop11_Lan1
\begin{ex}%[1D7H1-3]%[ma:34b1274f]
	Cho hàm số $y=\dfrac{1}{3} x^3+x^2-2x+1$ có đồ thị là $\left(C\right)$. Phương  trình tiếp tuyến của $\left(C\right)$ tại điểm $M\left(1;\dfrac{1}{3}\right)$ là
	\choice 
		{$y=3x-2$} 
		{$y=-3x+2$} 
		 {\True $y=x-\dfrac{2}{3}$}
		{$y=-x+\dfrac{2}{3}$} 
	 \loigiai{
		Ta có $y'=x^2+2x-2$ và  $y'(1)=1+2-2=1$.\\
		Phương trình  tiếp tuyến của $\left(C\right)$ tại điểm $M\left(1;\dfrac{1}{3} \right)$ là
		 $$y=y'(1)\left(x-1\right)+\dfrac{1}{3}=x-1+\dfrac{1}{3}=x-\dfrac{2}{3}.$$
	} 
\end{ex}


%%%%%%%%%%%%%%%% ĐÚNG/SAI %%%%%%%%%%%%%%%%

% nguồn: Ok_dungsai_lop11_Lan1
\begin{ex}%[1D7H1-3]%[ma:d543da4b]
 Cho hàm số $y=x^2+3 x+1$ có đồ thị $(C)$. Viết được phương trình tiếp tuyến của $(C)$ tại  giao điểm của $(C)$ với trục tung. Khi đó:
 Các mệnh  đề sau đúng hay sai?
  \choiceTF
  {\True Hệ số góc của phương trình tiếp tuyến bằng $ 3  $}
  {Phương trình tiếp  tuyến đi qua điểm $A(1; 3)$}
  {\True  Phương trình tiếp tuyến cắt đường thẳng $y=2 x+1$ tại điểm có hoành độ bằng $ 0 $}
  {\True Phương trình tiếp tuyến vuông góc với đường  thẳng $y=-\dfrac{1}{3} x+1$}
 \loigiai{ 
   Với $x_0=0 \Rightarrow y_0=1$.
  \\ 
  Ta  có
  $$ 
   f'(0)=\lim\limits_{x \rightarrow 0} \dfrac{f(x)-f(0)}{x-0}
  =\lim\limits_{x \rightarrow  0} \dfrac{x^2+3 x}{x}
  =\lim\limits_{x  \rightarrow 0}(x+3)=3.
  $$ 
  Vậy phương  trình tiếp tuyến là $y=3 x+1$.
  \\ 
  Xét tiếp tuyến $y=3 x+1$ và đường thẳng $y=-\dfrac{1}{3} x+1$ có tích hệ số góc $ 3 \cdot \dfrac{-1}{3}  = -1 $ nên hai đường thẳng này vuông góc với nhau.
 } 
\end{ex}

% nguồn: Ok_dungsai_lop11_Lan1
\begin{ex}%[1D7H1-3]%[ma:92a13155]
 Cho hàm số $f(x)=\dfrac{2}{1-x}$ có đồ thị $(C)$ và điểm $M(3;-1) \in(C)$. Khi đó, các mệnh  đề sau đúng hay sai?
  \choiceTF
  {\True Hệ số góc của  tiếp tuyến của $(C)$ tại điểm $M$ bằng $\dfrac{1}{2}$}
  {Phương trình tiếp tuyến của $(C)$ tại  $M$ song song với đường thẳng $y=-\dfrac{1}{2} x-\dfrac{5}{2}$}
  {\True Phương trình tiếp tuyến của $(C)$ tại $M$ vuông với đường thẳng $y=-2  x-\dfrac{5}{2}$}
  {\True Phương trình tiếp tuyến của $(C)$ tại $M$  đi qua điểm $A\left(0;-\dfrac{5}{2}\right)$}
 \loigiai{ 
  Ta có  $f'(x)=\left(\dfrac{2}{1-x}\right)'=\dfrac{2}{(1-x)^2}$ nên tiếp tuyến của $(C)$ tại $M$ có hệ số góc là
   $$
   f'(3)=\dfrac{2}{(1-3)^2}=\dfrac{1}{2}.
  $$ 
  Phương trình tiếp tuyến của $(C)$  tại $M$ là
  \begin{align*}
   y-(-1)&=\dfrac{1}{2}(x-3) \\ 
   y&=\dfrac{1}{2} x-\dfrac{5}{2}. 
  \end{align*}
  Thay tọa độ điểm $A\left(0;-\dfrac{5}{2}\right)$  vô phương trình tiếp tuyến trên, ta thấy
   $$
  -\dfrac{5}{2} = \dfrac{1}{2} \cdot 0 -\dfrac{5}{2}  \quad \text{đúng}.
   $$
  Vậy tiếp tuyến  của $ (C) $ tại $ M $ đi qua điểm $A\left(0;-\dfrac{5}{2}\right)$.
 } 
\end{ex}

% nguồn: Ok_dungsai_lop11_Lan1
\begin{ex}%[1D7H1-3]%[ma:80040eaa]
 Cho hàm số $y=f(x)=2 x^3$ có đồ thị  $(C)$ và điểm $M$ thuộc $(C)$ có hoành độ $x_0=-1$. Khi đó, các mệnh đề sau đúng hay sai?
 \choiceTF 
  {\True Hệ số góc của tiếp  tuyến của $(C)$ tại điểm $M$ bằng $ 6 $}
  {\True Phương trình tiếp tuyến của  $(C)$ tại $M$ đi qua điểm $A(0; 4)$}
  {Tiếp tuyến của $(C)$ tại $M$ cắt đường thẳng $d \colon y=3 x$ tại  điểm có hoành độ bằng $ 4 $}
  {\True Tiếp tuyến  của $(C)$ tại $M$ vuông góc với đường thẳng $d \colon y=-\dfrac{1}{6} x$}
 \loigiai{ 
  Ta có $f'(x)=\left(2 x^3\right)'=6 x^2$ nên tiếp tuyến của $(C)$ tại  $M$ có hệ số góc là
  $$ 
   f'(-1)=6 \cdot (-1)^2=6.
  $$ 
  Phương trình tiếp tuyến $ \Delta  $ của $(C)$ tại $M$ là
  \begin{align*}
   y-f(-1)&=6(x+1) \\ 
   y+2&=6(x+1)  \\
   \Delta \colon y&=6 x+4 . 
  \end{align*}
  Thay $ x=0 $ vào phương trình $ \Delta \colon y=6 x+4 $, ta được $ y=4 $. Do đó  tiếp tuyến $ \Delta $ đi qua điểm $ A(0;4) $.
  \\ 
  Xét phương  trình hoành độ giao điểm của tiếp tuyến $ \Delta $ và đường thẳng $d \colon y=3 x$
  $$ 
  6 x+4=3x 
   \Leftrightarrow
  x=-\dfrac{4}{3}  \ne 4.
   $$
   Xét tiếp tuyến $ \Delta \colon y=6x+4 $ và đường thẳng $d \colon y=-\dfrac{1}{6} x$, ta thấy
   $$
  k_{\Delta} \cdot k_{d} = 6 \cdot  \left(-\dfrac{1}{6}\right) = -1
   $$
  Suy ra $ \Delta  \perp d $.
  }
\end{ex}

% nguồn: Ok_dungsai_lop11_Lan1
\begin{ex}%[1D7H1-3]%[ma:ede584f8]
 Viết được phương trình tiếp tuyến của đồ thị hàm số $y=\dfrac{x+9}{x+1}$ biết  tiếp tuyến vuông góc với đường thẳng $d: x-2 y+2=0$. Khi đó
 \choiceTF 
   {\True Có $ 2 $ phương trình tiếp tuyến thỏa mãn.}
  {\True Hệ số góc của tiếp tuyến bằng  -2}
  {\True Tiếp tuyến đi qua điểm  $A(1; 5)$}
  {\True Tiếp tuyến  đi qua điểm $B(-1;-7)$}
 \loigiai{ 
  Đường thẳng $d: x-2 y+2=0 \Rightarrow y=\dfrac{1}{2} x+1$ nên đường thẳng $d$ có hệ số  góc là $k_d=\dfrac{1}{2}$.
  Tiếp tuyến cần tìm có  hệ số góc $k$ vuông góc với đường thẳng $d$
  $$ 
   \Rightarrow k \cdot k_d=-1 \Rightarrow k=-\dfrac{1}{k_d}=-2 \text {. }
  $$ 
  Hoành độ tiếp điểm là nghiệm của  phương trình
   $$
  y'=k 
  \Rightarrow 
  \dfrac{-8}{(x+1)^2}=-2 
  \Rightarrow 
   \heva{&x=1 \\ &x=-3}
   $$
  Với $x=1$, phương trình tiếp tuyến  là $d_1 \colon y=-2 x+7$.
  \\ 
  Với $x=-3$,  phương trình tiếp tuyến là $d_2 \colon y=-2 x-9$.
  \\ 
  Thay $ x=1 $ vào phương trình tiếp  tuyến $d_1 \colon y=-2 x+7$, ta được $ y=5 $. Suy ra $ d_1 $ đi qua $ A $.
   \\
  Thay $ x=-1 $ vào phương trình tiếp tuyến $d_2 \colon y=-2  x-9$, ta được $ y=-7 $. Suy ra $ d_2 $ đi qua $ B $.
  }
\end{ex}

% nguồn: Ok_dungsai_lop11_Lan2
\begin{ex}%[1D7V1-3]%[ma:4b9d8abd]
	Cho hàm số $y=x^2+2x-4$ có đồ thị  $(C)$.
	 \choiceTF
		{Hệ số góc  của tiếp tuyến của $(C)$ tại điểm có hoành độ $x_0=1$ thuộc $(C)$ là $-4$}
		{\True Có hai phương trình tiếp tuyến của đồ thị hàm số tại  điểm có tung độ $y_0=-1$ thuộc $(C)$}
		{Phương trình tiếp tuyến của đồ thị hàm số  biết hệ số góc tiếp tuyến bằng $-4$ là $y=-4x-12$}
		{\True Phương trình tiếp tuyến của đồ thị hàm số biết tiếp tuyến đó song song với đường thẳng $y=1-3x$ là $y=-3x-\dfrac{41}{4}$} 
	\loigiai{ 
		Tại $x_0 \in  \mathbb{R}$ tùy ý, ta có
		\begin{align*}
			f(x)-f\left(x_0\right)&=x^2+2x-4-\left(x_0  ^2+2x_0-4\right)\\
			&=x^2-x_0 ^2+2x-2x_0\\ 
			 &=\left(x-x_0\right)\left(x+x_0+2\right).
		\end{align*}
		\begin{align*}
			\lim\limits_{x\to  x_0}\dfrac{f(x)-f\left(x_0\right)}{x-x_0}&=\lim\limits_{x\to x_0}\dfrac{\left(x-x_0\right)\left(x+x_0+2\right)}{x-x_0}\\
			&=\lim\limits_{x\to x_0}\left(x+x_0+2\right)\\ 
			&=2x_0+2. 
		\end{align*}
		 Ta có $y'=2x+2$.
		\begin{itemchoice}
			 \itemch Hệ số góc của tiếp tuyến của $(C)$ tại điểm có hoành độ $x_0=1$ thuộc $(C)$ là $k=y'(1)=4$.
			\itemch  Ta có $y_0=-1\Rightarrow x_0^2+2x_0-4=-1\Leftrightarrow \hoac{&x_0=1\Rightarrow y_0=-1\\&x_0=-3\Rightarrow y_0=-1.}$
			\begin{itemize}
				 \item Phương trình tiếp tuyến tại điểm $(1;-1)$ là $y=y'(1)\cdot(x-1)+y(1)\Leftrightarrow y=4x-5$.
				\item Phương trình tiếp tuyến tại điểm $(-3;-1)$ là  $y=y'(-3)\cdot(x+3)+y(-3)\Leftrightarrow y=-4x-13$.
			\end{itemize}
			\itemch Gọi $M\left(x_0;y_0\right)$ là tiếp điểm của  tiếp tuyến của đồ thị $(C)$ với hệ số góc $k=-4$.\\
			 Ta có $y'(x_0)=-4\Leftrightarrow 2x_0+2=-4\Leftrightarrow x_0=-3\Rightarrow y_0=-1$.\\
			Phương trình tiếp tuyến  với hệ số góc $k=-4$ là $y=-4(x+3)-1\Leftrightarrow y=-4x-13$.
			\itemch Vì tiếp tuyến song song với đưởng  thẳng $y=1-3x$ nên tiếp tuyến có hệ số góc $k=-3$.\\
			Gọi $M\left(x_0;y_0\right)$ là tiếp điểm của  tiếp tuyến của đồ thị $(C)$ với hệ số góc $k=-4$.\\
			Ta có $y' (x_0)=-3\Leftrightarrow 2x_0+2=-3\Leftrightarrow x_0=-\dfrac{5}{2}  \Rightarrow y_0=-\dfrac{11}{4}$.\\
			Phương trình tiếp tuyến với hệ số góc $k=-3$ là $y=-3\Bigl(x+\dfrac{5}{2} \Bigr)-\dfrac{11}{4} \Leftrightarrow  y=-3x-\dfrac{41}{4}$.
		\end{itemchoice}
	 }
\end{ex}


%%%%%%%%%%%%%%%% TRẢ LỜI NGẮN %%%%%%%%%%%%%%%%

% nguồn: Ok_TraLoiNgan_Lop11_Lan1
\begin{ex}%[1D7H1-3]%[ma:4a1513d3]
	Tiếp tuyến với đồ thị hàm số $y=x^3+2 x^2-3$. Tại điểm có hoành  độ bằng $ -2 $ có hệ số góc là bao nhiêu?
	\shortans{4} 
	 \loigiai{
		Ta  có: $x_0=-2 \Rightarrow y_0=(-2)^3+2 \cdot(-2)^2-3=-3$
		 $$
		f'(-2)=\lim\limits_{x \rightarrow-2} \dfrac{f(x)-f(2)}{x+2} 
		=\lim\limits_{x \rightarrow-2} \dfrac{x^3+2 x^2}{x+2} 
		=\lim\limits_{x  \rightarrow-2} \dfrac{x^2(x+2)}{x+2}
		 =\lim\limits_{x \rightarrow-2} x^2
		=4. 
		 $$
		Vậy hệ số góc của tiếp tuyến cần tìm là $ 4 $. 
	} 
\end{ex}

% nguồn: Ok_TraLoiNgan_Lop11_Lan1
\begin{ex}%[1D7H1-3]%[ma:4943769a]
	Cho hàm số $y=f(x)=\dfrac{x+1}{3 x}$ có đồ thị $(C)$. 
	Phương trình tiếp tuyến của đồ thị hàm số $(C)$ tại giao  điểm của $(C)$ với trục hoành có dạng $ y=ax+b $. Khi đó, giá trị của biểu thức $ a+b $ là bao nhiêu? (làm tròn đến hàng phần mười)
	 \shortans{-0,7}
	\loigiai{ 
		Toạ độ giao điểm của $(C)$ với  trục hoành là điểm $(-1; 0)$.
		Phương trình tiếp tuyến của $(C)$ tại điểm $(-1; 0)$  là
		\begin{align*}
			y-0&=f'(-1)(x+1)\\ 
			 y&=-\dfrac{1}{3}(x+1)\\
			y&=-\dfrac{1}{3} x-\dfrac{1}{3}. 
		\end{align*}
		Do đó $ a=-\dfrac{1}{3} $ và $ b=-\dfrac{1}{3} $. Suy ra $ a+b  = -\dfrac{2}{3} \approx -0{,}7 $.
	 }
\end{ex}

% nguồn: Ok_TraLoiNgan_Lop11_Lan1
\begin{ex}%[1D7H1-3]%[ma:5c07d113]
	Cho hàm số $y=f(x)=-2  x^3+x$ có đồ thị $(C)$.
	Phương trình tiếp tuyến của đồ thị $(C)$ tại điểm $M(-2; 14)$ có dạng $ y=mx+n $. Khi đó, giá trị  của $ m+n $ là bao nhiêu?
	 \shortans{-55}
	\loigiai{ 
		Ta có 
		$f'(x)=\left(-2 x^3+x\right)'=-6  x^2+1$
		nên hệ số góc của tiếp tuyến của $(C)$  tại điểm $M(-2; 14)$ là
		$f'(-2)=-6 \cdot(-2)^2+1=-23$. 
		\\ 
		Phương trình tiếp  tuyến của $(C)$ tại điểm $M$ là:
		$$ 
		 y-14=-23(x+2) \Leftrightarrow y=-23 x-32.
		$$ 
		Do đó $ m=-23 $, $ n=-32 $  và $ m+n=-55 $.
	 }
\end{ex}

% nguồn: Ok_TraLoiNgan_Lop11_Lan1
\begin{ex}%[1D7H1-3]%[ma:fb3611f7]
	Cho  hàm số $y=f(x)=\dfrac{x+1}{3 x}$ có đồ thị $(C)$.
	Tìm được bao nhiêu phương trình tiếp tuyến của đồ thị  hàm số $(C)$ tại giao điểm của $(C)$ với đường thẳng $y=x+1$?
	\shortans{2} 
	 \loigiai{
		Toạ độ giao điểm của  $(C)$ với đường thẳng $y=x+1$ là nghiệm của hệ phương trình
		 $$
		\heva{&y=\dfrac{x+1}{3 x}  \\ &y=x+1}
		\Leftrightarrow 
		 \heva{&\dfrac{x+1}{3 x}=x+1 \\ &y=x+1}
		\Leftrightarrow 
		\heva{&3 x^2+2 x-1=0  &&(1) \\ &y=x+1. &&(2)}
		$$ 
		Phương trình $  (1) $ có hai nghiệm là $ x=-1 $ và $ x=\dfrac{1}{3} $.
		 \\
		Vậy ta tìm được $ 2  $ tiếp tuyến.
	} 
\end{ex}

% nguồn: Ok_TraLoiNgan_Lop11_Lan2
\begin{ex}%[1D7H1-3]%[ma:ca078a4a]
	Tính đạo hàm của hàm số $f(x)=\heva{&\dfrac{\sqrt{x^3+x^2+1}-1}{x}\text{ khi }x\ne 0\\&0\text{  khi }x=0}$ tại $x=0$.
	\shortans{0,5} 
	\loigiai{ 
		 Ta có $f(0)=0$.
		\begin{eqnarray*}
			\lim\limits_{x\to 0}\dfrac{f(x)-f(0)}{x}&=&\lim\limits_{x\to  0}\dfrac{\sqrt{x^3+x^2+1}-1}{x^2}\\
			&=&\lim\limits_{x\to  0}\dfrac{x+1}{\sqrt{x^3+x^2+1}+1}\\
			 &=&\dfrac{1}{2}.
		\end{eqnarray*}
		 Vậy $f'(0)=\dfrac{1}{2}$.
	} 
\end{ex}

% nguồn: Ok_TraLoiNgan_Lop11_Lan2
\begin{ex}%[1D7H1-3]%[ma:ab0c3f7c]
	Cho hàm số $f(x)=\dfrac{1}{x+2}$ có đồ thị $(C)$. Tìm hệ số góc của tiếp tuyến với đồ thị (C) tại  điểm có hoành độ bằng $-3$.
	 \shortans{-1}
	 \loigiai{
		Hệ số góc của tiếp tuyến với đồ thị $(C)$  tại $x_0=-3$ là
		\[f'(-3)=\lim\limits_{x\to-3} \dfrac{f(x)-f(-3)}{x+3}=\lim\limits_{x\to-3}  \dfrac{\tfrac{1}{x+2}+1}{x+3}=\lim\limits_{x\to-3}\dfrac{1}{x+2}=-1.
		\] 
	 }
\end{ex}

% nguồn: Ok_TraLoiNgan_Lop11_Lan1
\begin{ex}%[1D7V1-3]%[ma:2c4193fc]
	Cho hàm số  $y=f(x)=-2 x^3+x$ có đồ thị $(C)$.
	Tính hệ số góc của tiếp tuyến của đồ thị $(C)$ tại điểm có hoành độ  bằng $ 1 $.
	 \shortans{-5}
	\loigiai{ 
		Ta có $f'(x)=\left(-2 x^3+x\right)'=-6 x^2+1$  nên hệ số góc của tiếp tuyến của $(C)$ tại điểm có hoành độ bằng $ 1 $ là $f'(1)=-6 \cdot(1)^2+1=-5$.
	} 
\end{ex}

% nguồn: Ok_TraLoiNgan_Lop11_Lan2
\begin{ex}%[1D7V1-3]%[ma:374ad5cd]
	Cho hàm số $y=\dfrac{2x-1}{x-2}$ có đồ thị $(C)$. Phương trình tiếp tuyến của đồ thị $(C)$, biết tung độ tiếp điểm bằng $3$ có dạng $y=ax+b$.  Tính $6a+3b$.
	 \shortans{12}
	 \loigiai{
		Tại điểm  có tung độ tiếp điểm bằng $3$, ta có
		\begin{itemize}
			\item  $\dfrac{2x-1}{x-2}=3\Leftrightarrow x=5$.
			\item $f'(5)=\lim\limits_{x\to  5} \dfrac{f(x)-f(5)}{x-5}=\lim\limits_{x\to 5} \dfrac{\tfrac{2x-1}{x-2}-3}{x-5}=\lim\limits_{x\to 5} \dfrac{-(x-5)}{(x-5)(x-2)}=\lim\limits_{x\to 5} \dfrac{-1}{x-2}=-\dfrac{1}{3}$.
		\end{itemize}
		Phương trình tiếp tuyến cần  tìm là $y=-\dfrac{1}{3} (x-5)+3\Leftrightarrow y=-\dfrac{1}{3} x+\dfrac{14}{3}$.\\
		Suy ra $a=-\dfrac{1}{3}$; $b=\dfrac{14}{3}$.  Do đó $6a+3b=12$.
	} 
\end{ex}

% nguồn: Ok_TraLoiNgan_Lop11_Lan2
\begin{ex}%[1D7V1-3]%[ma:e7b98411]
	Cho hàm số $f(x)=\dfrac{1}{x}$ có đồ thị $(C)$ và điểm $M(1;1)\in (C)$. Phương trình tiếp tuyến với $(C)$  tại điểm $M$ có dạng $y=ax+b$. Tính $T=2a+b$.
	 \shortans{0}
	\loigiai{ 
		Hệ sộ góc của tiếp tuyến tại  điểm $M\left(1;1\right)\in (C)$ là
		 \[k=f'(1)=\lim\limits_{x\to 1}\dfrac{f(x)-f(1)}{x-1}=\lim\limits_{x\to 1}\dfrac{\tfrac{1}{x}-1}{x-1}=\lim\limits_{x\to 1}\dfrac{-1}{x}=-1.\]
		Phương trình tiếp tuyến $M(1;1)\in (C)$ có dạng  $y=f'(1)(x-1)+1\Leftrightarrow y=-x+2\Rightarrow \heva{&a=-1\\&b=2.}$\\
		 Vậy $T=2a+b=2\cdot (-1)+2=0$.
	} 
\end{ex}

% nguồn: Ok_TraLoiNgan_Lop11_Lan2
\begin{ex}%[1D7V1-3]%[ma:4b8b1976]
	Cho hàm số $f(x)=\dfrac{x-1}{x+1}$ có đồ thị $\left(C\right)$. Phương  trình tiếp tuyến với $\left(C\right)$ vuông góc với đường thẳng $\Delta\colon y=-\dfrac{1}{2} x+3$ có dạng $y=ax+b$. Tính $T=a^2\cdot b$. Biết tọa độ tiếp điểm có hoành độ âm.
	 \shortans{28}
	\loigiai{ 
		Ta có \begin{eqnarray*}
		f'(x)&=&\lim\limits_{h\to  0}\dfrac{f\left(x+h\right)-f(x)}{h}\\
		&=&\lim\limits_{h\to  0}\dfrac{\dfrac{\left(x+h\right)-1}{\left(x+h\right)+1}-\dfrac{x-1}{x+1}}{h}\\
		&=&\lim\limits_{h\to  0}\dfrac{2}{\left(x+1\right)\left(x+h+1\right)}\\
		 &=&\dfrac{2}{\left(x+1\right)^2}.
	\end{eqnarray*}
	 Suy ra $f'(x)=\dfrac{2}{\left(x+1\right)^2}$.\\
	 Gọi $M\left(x_0;\dfrac{x_0-1}{x_0+1} \right)\in \left(C\right)$ là tọa độ tiếp điểm.\\
	Hệ  số góc của tiếp tuyến tại điểm $M\in \left(C\right)$ là $k=f'(x_0)=\dfrac{2}{\left(x_0+1\right)^2}$.\\
	Tiếp  tuyến tại điểm $M\in \left(C\right)$ vuông góc với đường thẳng $\Delta\colon y=-\dfrac{1}{2} x+3$ nên
	 \allowdisplaybreaks
	\begin{eqnarray*}
		 &&k\cdot\left(-\dfrac{1}{2} \right)=-1\\
		&\Leftrightarrow& \dfrac{1}{\left(x_0+1\right)^2}=1\Leftrightarrow \hoac{&x_0=0\text{ (loại do tọa độ tiếp điểm có  hoành độ âm)}\\&x_0=-2.}
	\end{eqnarray*}
	Với $x_0=-2\Rightarrow M\left(-2;3\right)$. Khi đó , phương trình tiếp tuyến tại điểm $M(-2;3)\in (C)$ có dạng  \[y=2\left(x+2\right)+3\Leftrightarrow y=2x+7\Rightarrow \heva{&a=2\\&b=7.}\]
	Vậy  $T=a^2\cdot b=28$.
 }
\end{ex}

% nguồn: Ok_TraLoiNgan_Lop11_Lan2
\begin{ex}%[1D7V1-3]%[ma:01acc5f4]
	Phương trình tiếp tuyến với đồ thị  hàm số $y=x^3+2x^2-3$ tại điểm có hoành độ bằng $-2$ có dạng $y=ax+b$. Tính $2a+3b$.
	 \shortans{23}
	\loigiai{ 
		Ta có 
		\begin{itemize}
			\item $x_0=-2\Rightarrow y_0=(-2)^3+2\cdot (-2)^2-3=-3$. 
			\item $f'(-2)=\lim\limits_{x\to-2} \dfrac{f(x)-f(2)}{x+2}=\lim\limits_{x\to-2} \dfrac{x^3+2x^2}{x+2}=\lim\limits_{x\to-2} \dfrac{x^2(x+2)}{x+2}=\lim\limits_{x\to-2}  x^2=4$.
		\end{itemize}
		Khi đó phương trình tiếp tuyến  là $y=y' \left(x_0 \right)\left(x-x_0 \right)+y_0=4(x+2)-3=4x+5$.\\
		Suy ra $a=4$; $b=5$. Do đó  $2a+3b=23$.
	 }
\end{ex}

% nguồn: Ok_TraLoiNgan_Lop11_Lan2
\begin{ex}%[1D7V1-3]%[ma:4f65636b]
	Cho hàm  số $y=\dfrac{x+1}{3x}$. Tính hệ số góc tiếp tuyến của đồ thị hàm số tại điểm có hoành độ bằng $3$. (Kết quả làm tròn đến hàng phần trăm).
	\shortans{0,04} 
	 \loigiai{
		Tại $x_0 \in \mathbb{R}\setminus \{0\}$  tùy ý, ta có
		\begin{itemize}
			\item $f(x)-f\left(x_0 \right)=\dfrac{x+1}{3x}-\dfrac{x_0+1}{3x_0}=\dfrac{-\left(x-x_0 \right)}{3x\cdot  x_0}$.
			\item  $\dfrac{f(x)-f\left(x_0 \right)}{x-x_0}=\dfrac{-\left(x-x_0 \right)}{3x\cdot x_0}\cdot \dfrac{1}{x-x_0}=\dfrac{-1}{3x\cdot x_0}$.
			\item $\lim\limits_{x\to x_0}\dfrac{f(x)-f\left(x_0 \right)}{x-x_0}=\lim\limits_{x\to x_0}\dfrac{-1}{3x\cdot x_0}=\dfrac{-1}{3x_0 ^2}$. 
		\end{itemize}
		Suy ra $y'=-\dfrac{1}{3x^2}$.\\ 
		Khi đó, hệ số góc tiếp tuyến của đồ thị hàm số tại điểm có hoành độ bằng  $3$ là \[y'(3)=-\dfrac{1}{27} \approx 0{,}04.\]
	 }
\end{ex}


%%%%%%%%%%%%%%%% TỰ LUẬN %%%%%%%%%%%%%%%%

% nguồn: Ok_TuanLuan_Toan11_Lan1
\begin{bt}%[1D7H1-3]%[ma:9f4969e0]
	Viết phương trình tiếp tuyến với đồ thị hàm số $y=x^3+2x^2-3$ tại  điểm có hoành độ bằng $-2$.
	\loigiai{ 
		 Tại $x_0=-2$, ta có
		\begin{itemize}
			\item  $y_0=(-2)^3+2\cdot (-2)^2-3=-3$.
			\item $f'(-2)=\lim\limits_{x\to -2}\dfrac{f(x)-f(-2)}{x+2}=\lim\limits_{x\to -2}\dfrac{x^3+2x^2}{x+2}=\lim\limits_{x\to  -2}\dfrac{x^2(x+2)}{x+2}=\lim\limits_{x\to -2} x^2=4$.
		\end{itemize}
		Vậy phương trình tiếp tuyến  là $y=y'(x_0)\left(x-x_0\right)+y_0=4(x+2)-3=4x+5$.
	 }
\end{bt}

% nguồn: Ok_TuanLuan_Toan11_Lan1
\begin{bt}%[1D7H1-3]%[ma:bb6462b7]
	Cho hàm số $f(x)=\dfrac{1}{x+2}$ có đồ thị $(C)$. Tìm hệ số góc của tiếp tuyến với đồ thị $(C)$  tại điểm có hoành độ bằng $-3$.
	\loigiai{  \\
		Hệ số góc của tiếp tuyến với  đồ thị $(C)$ tại $x_0=-3$ là \\
		$f'(-3)=\lim\limits_{x \to -3} \dfrac{f(x)-f(-3)}{x+3}=\lim\limits_{x \to -3} \dfrac{\frac{1}{x+2}+1}{x+3}=\lim\limits_{x \to -3}  \dfrac{1}{x+2}=-1$.\\
	} 
\end{bt}

% nguồn: Ok_TuanLuan_Toan11_Lan2
\begin{ex}%[1D7H1-3]%[ma:823c28e6]
	Viết phương trình tiếp  tuyến với đồ thị hàm số $y=x^3+2x^2-3$ tại điểm có hoành độ bằng $-2$.
	\loigiai{ 
		Tại  $x_0=-2$, ta có
		\begin{itemize}
			\item $y_0=(-2)^3+2\cdot  (-2)^2-3=-3$.
			 \item $f'(-2)=\lim\limits_{x\to -2}\dfrac{f(x)-f(-2)}{x+2}=\lim\limits_{x\to -2}\dfrac{x^3+2x^2}{x+2}=\lim\limits_{x\to -2}\dfrac{x^2(x+2)}{x+2}=\lim\limits_{x\to -2} x^2=4$.
		\end{itemize}
		Vậy phương trình tiếp  tuyến là $y=y'(x_0)\left(x-x_0\right)+y_0=4(x+2)-3=4x+5$.
	} 
\end{ex}

% nguồn: Ok_TuanLuan_Toan11_Lan2
\begin{ex}%[1D7H1-3]%[ma:5dfd46d3]
	Cho hàm số $f(x)=\dfrac{1}{x+2}$ có đồ thị $(C)$. Tìm hệ số góc của tiếp tuyến với đồ  thị $(C)$ tại điểm có hoành độ bằng $-3$.
	\loigiai{ \\ 
		Hệ số góc của  tiếp tuyến với đồ thị $(C)$ tại $x_0=-3$ là \\
		 $f'(-3)=\lim\limits_{x \to -3} \dfrac{f(x)-f(-3)}{x+3}=\lim\limits_{x \to -3} \dfrac{\frac{1}{x+2}+1}{x+3}=\lim\limits_{x \to -3} \dfrac{1}{x+2}=-1$.\\
	 }
\end{ex}

% nguồn: Ok_TuanLuan_Toan11_Lan2
\begin{ex}%[1D7H1-3]%[ma:03cee886]
	 Cho hàm số $ y=x^{2}+2 x-4 $ có đồ thị $ (C) $.
	\begin{enumerate}
		\item  Tìm hệ số góc của tiếp tuyến của $ (C) $ tại điểm có hoành độ $ x_{0}=1 $ thuộc $ (C) $.
		\item Viết phương trình tiếp tuyến của đồ thị hàm số tại điểm có hoành độ $ x_{0}=0 $ thuộc  $ (C) $.
		\item Viết  phương trình tiếp tuyến của đồ thị hàm số tại điểm có tung độ $ y_{0}=-1 $ thuộc $ (C) $.
		\item Viết phương trình tiếp tuyến của đồ thị hàm số biết hệ số góc của tiếp tuyến bằng  $-4$.
		\item Viết phương trình tiếp tuyến của đồ thị hàm số, biết tiếp  tuyến đó song song với đường thẳng $ y=1-3 x $.
	\end{enumerate}
	 \loigiai{
		Tại $ x_{0} \in \mathbb{R}  $ tùy ý, ta có
		\begin{eqnarray*}
			f(x)-f\left(x_{0}\right)&=&x^{2}+2 x-4-x_{0}^{2}-2 x_{0}+4\\ 
			&=&\left(x-x_{0}\right)\left(x+x_{0}+2\right) \\ 
			\dfrac{f(x)-f\left(x_{0}\right)}{x-x_{0}}&=&\dfrac{\left(x-x_{0}\right)\left(x+x_{0}+2\right)}{x-x_{0}}=x+x_{0}+2  \\
			\lim\limits_{x \rightarrow x_{0}} \dfrac{f(x)-f\left(x_{0}\right)}{x-x_{0}}&=&\lim\limits_{x \rightarrow x_{0}}\left(x+x_{0}+2\right)=2 x_{0}+2. 
		\end{eqnarray*}
		Suy ra $y'=2  x+2$.
		\begin{enumerate}
			\item Hệ số góc  của tiếp tuyến của $ (C) $ tại điểm có hoành độ $ x_{0}=1 $ thuộc $ (C) $ là $ k=y'(1)=4 $.
			\item Phương trình tiếp tuyến của đồ thị hàm số tại điểm có hoành độ $ x_{0}=0 $ thuộc (C)  là
			\[ 
			y=y'(0)(x-0)+y(0) \Leftrightarrow y=2 x-4. 
			\] 
			\item Với $ y_{0}=-1 \Rightarrow y=x_{0}^{2}+2 x_{0}-4=-1 \Leftrightarrow\hoac{&x_{0}=1  \\& x_{0}=-3.}$\\
			Vậy có hai tiếp điểm thuộc $ (C) $ có tung độ $ y_{0}=-1 $ là  $ (1;-1) $ và $ (-3;-1) $.\\
			Phương  trình tiếp tuyến tại điểm $ (1;-1) $ là $$ y=y'(1)(x-1)+y(1) \Leftrightarrow y=4 x-5.$$
			Phương trình tiếp tuyến tại điểm $ (-3;-1)  $ là $$ y=y'(-3)(x+3)+y(-3) \Leftrightarrow y=-4 x-13.$$
			\item Gọi $ M(a; b) $ là tiếp điểm của tiếp tuyến của đồ thị ( $ C $) với hệ số góc $ k=-4 $ nên 
			 \[
			 y'(a)=-4 \Leftrightarrow 2 a+2=-4 \Leftrightarrow a=-3.
			\] 
			Suy ra  $b=-1$.\\
			Vậy phương  trình tiếp tuyến với hệ số góc $ k=-4 $ là $$ y=-4(x+3)-1 \Leftrightarrow y=-4 x-13.$$
			\item Ta có $y'  = -3$.\\
			Xét $2x+2 = -3 \Rightarrow x = -\dfrac{5}{2}$ nên $y\left(-\dfrac{5}{2}\right) =  -\dfrac{11}{4}$.\\
			Phương trình tiếp tuyến là $$y + \dfrac{11}{4} = -3\left(x+\dfrac{5}{2}\right) \Leftrightarrow y = -3x -  \dfrac{37}{4}.$$
			\item Vì tiếp tuyến đó song song với đưởng thẳng $ y=1-3 x $ nên tiếp tuyến có  hệ số góc $ k=-3 $.\\
			Gọi $ M(a; b) $ là tiếp điểm của tiếp tuyến của đồ thị ( $ c $) với hệ số góc $ k=-4 $, ta có 
			\[ 
			y'(a)=-3  \Leftrightarrow 2 a+2=-3 \Leftrightarrow a=-\dfrac{5}{2}.
			\] 
			Suy ra $b=-\dfrac{11}{4}$.\\ 
			Vậy phương trình tiếp  tuyến với hệ số góc $ k=-3 $ là $$ y=-3\left(x+\dfrac{5}{2}\right)-\dfrac{11}{4} \Leftrightarrow y=-3 x-\dfrac{41}{4}.$$
		\end{enumerate}
	} 
\end{ex}

% nguồn: Ok_TuanLuan_Toan11_Lan1
\begin{bt}%[1D7V1-3]%[ma:3fbf7a36]
	Cho hàm số $y = x^3 + 2x^2+1$ có đồ thị  là $(C)$. Viết phương trình tiếp tuyến của $(C)$ tại điểm $M(1; 4)$.
	 \loigiai{
		Đặt $f(x) = x^3 + 2x^2 + 1 \Rightarrow f'(x) = 3x^2  + 4x$.
		\\Phương trình tiếp tuyến của $(C)$ tại $M(1;4)$  có dạng
		\\$ y = f'(1)  \cdot (x-1) + 4 \Rightarrow y = 7x -3$.
	 }
\end{bt}

% nguồn: Ok_TuanLuan_Toan11_Lan2
\begin{ex}%[1D7V1-3]%[ma:3aeadaf8]
	Cho hàm số $y = x^3  + 2x^2+1$ có đồ thị là $(C)$. Viết phương trình tiếp tuyến của $(C)$ tại điểm $M(1; 4)$.
	\loigiai{ 
		Đặt $f(x) = x^3 + 2x^2  + 1 \Rightarrow f'(x) = 3x^2 + 4x$.
		\\Phương trình tiếp tuyến của $(C)$ tại  $M(1;4)$ có dạng
		\\$ y  = f'(1) \cdot (x-1) + 4 \Rightarrow y = 7x -3$.
	} 
\end{ex}

% nguồn: Ok_TuanLuan_Toan11_Lan2
\begin{ex}%[1D7V1-3]%[ma:f844508e]
	Cho hàm số $ y=\dfrac{x+1}{3 x} $ có đồ  thị $(C)$.
	\begin{enumerate}
		\item Viết  phương trình tiếp tuyến của đồ thị hàm số tại giao điểm của $ (C) $ với trục $ O y $.
		\item Viết phương trình tiếp tuyến của đồ thị hàm số tại giao  điểm của $ (C) $ với trục $ O x $.
		\item Viết phương trình tiếp tuyến của đồ thị hàm số tại giao  điểm của $ (C) $ với đường thẳng $ y=x+1 $.
		\item Viết phương trình tiếp tuyến của đồ thị  hàm số biết hệ số góc của tiếp tuyến bằng $ k=-\dfrac{1}{3} $.
		\item Viết phương trình tiếp tuyến của đồ thị hàm số, biết  tiếp tuyến đó vuông góc với đưởng thẳng $ y=3 x-4 $.
	\end{enumerate}
	\loigiai{ 
		Tại $  x_{0} \in \mathbb{R} \backslash\{0\} $ tùy ý, ta có
		\begin{eqnarray*}
			f(x)-f\left(x_{0}\right)&=&\dfrac{x+1}{3 x}-\dfrac{x_{0}+1}{3 x_{0}}=\dfrac{x_{0}(x+1)-x(x_{0}+1)}{3 x x_{0}} \\ 
			&=& \dfrac{x_{0}x+x_{0}-x  x_{0}-x}{3 x x_{0}} = \dfrac{x_{0}-x}{3 x x_{0}} = \dfrac{-(x-x_{0})}{3 x x_{0}} \\
			\dfrac{f(x)-f\left(x_{0}\right)}{x-x_{0}}&=&\dfrac{-(x-x_{0})}{3 x x_{0}}  \cdot \dfrac{1}{x-x_{0}}=\dfrac{-1}{3 x x_{0}} \\
			\lim\limits_{x  \rightarrow x_{0}} \dfrac{f(x)-f\left(x_{0}\right)}{x-x_{0}}&=&\lim\limits_{x \rightarrow x_{0}} \dfrac{-1}{3 x x_{0}}=\dfrac{-1}{3 x_{0}^{2}}.
		\end{eqnarray*}
		 Suy ra $y'=-\dfrac{1}{3 x^{2}}$.
		\begin{enumerate}
			\item Vì $ (C) $ không cắt $ O y $ nên không tồn tại tiếp tuyến thỏa yêu  cầu bài toán.
			\item Tọa độ giao điểm của $ (C)  $ với trục $ O x $ là $ (-1; 0) $.\\
			Suy ra phương trình tiếp tuyến tại giao điểm của $ (C) $  với trục $ O x $ là
			\[ 
			y=y'\left(-1\right)(x+1)+0 \Leftrightarrow y=-\dfrac{1}{3} x-\dfrac{1}{3}. 
			 \]
			\item Tọa độ giao điểm của $ (C) $ với đường thẳng  $ y=x+1 $ là nghiệm của phương trình
			 \[
			 \dfrac{x+1}{3 x}=x+1 \Rightarrow\hoac{&x=-1\\&x=\dfrac{1}{3}.}
			\] 
			\begin{itemize}
				\item Tại $x=-1$  nên $y=0$.\\
				Phương trình tiếp tuyến tại điểm $ (-1; 0) $ là  $$ y=y'\left(-1\right)(x+1)+0 \Leftrightarrow y=-\dfrac{1}{3} x-\dfrac{1}{3}.$$
				\item Tại $x=\dfrac{1}{3}$ nên  $y=\dfrac{4}{3}$.\\
				Phương trình tiếp tuyến tại điểm $ \left(\dfrac{1}{3}; \dfrac{4}{3}\right) $ là $$ y=y'\left(\dfrac{1}{3}\right)\left(x-\dfrac{1}{3}\right)+\dfrac{4}{3}  \Leftrightarrow y=-3\left(x-\dfrac{1}{3}\right)+\dfrac{4}{3} \Leftrightarrow y=-3x+3.$$
			\end{itemize}
			\item Gọi $ M(a; b)  $ là tiếp điểm của tiếp tuyến của đồ thị $ (C) $ với hệ số góc $ k=-\dfrac{1}{3} $ nên
			\[ 
			 y'(a)=-\dfrac{1}{3} \Leftrightarrow-\dfrac{1}{3 a^{2}}=-\dfrac{1}{3} \Leftrightarrow\hoac{&a=1\\&a=-1.}
			\] 
			\begin{itemize}
				 \item Khi $a=1$ thì $b=\dfrac{2}{3}$.\\
				Phương trình  tiếp tuyến với hệ số góc $ k=-\dfrac{1}{3} $ là $$ y=-\dfrac{1}{3}(x-1)+\dfrac{2}{3} \Leftrightarrow y=-\dfrac{1}{3} x+1.$$
				\item  Khi $a=-1$ thì $b=0$.\\
				Phương trình tiếp tuyến với hệ số góc  $ k=-\dfrac{1}{3} $ là $y=-\dfrac{1}{3} x-\dfrac{1}{3}$.
			\end{itemize}
			\item Tiếp tuyến đó  vuông góc với đưởng thẳng $ y=3 x-4 $.\\
			Suy ra tiếp tuyến hệ số góc $  k=-\dfrac{1}{3} $.\\
			Ta biết đạo hàm của hàm số là $y'=-\dfrac{1}{3x^{2}}$. Hệ số góc của tiếp tuyến tại điểm có hoành độ $x_{0}$ chính là  $y'(x_{0})$, ta có
			\[ 
			y'(x_{0})=-\dfrac{1}{3}\Leftrightarrow -\dfrac{1}{3x_{0}^{2}}=-\dfrac{1}{3}\Leftrightarrow  x_0=\pm 1.
			 \]
			\begin{itemize}
				\item Khi $x_0=1$ thì  $y_0=\dfrac{2}{3}$.\\
				Phương trình  tiếp tuyến (với hệ số góc $k=-\dfrac{1}{3}$) là
				 \[
				y-\dfrac{2}{3}=-\dfrac{1}{3}\left(x-1\right)\Leftrightarrow y=-\dfrac{1}{3}x+1. 
				\] 
				 \item Khi $x_0=-1$ thì $y_0=0$.\\
				Phương trình tiếp tuyến  là
				 \[
				 y-0=-\dfrac{1}{3}\left(x+1\right)\Leftrightarrow y=-\dfrac{1}{3}x-\dfrac{1}{3}.
				 \]
			\end{itemize}
			Vậy các tiếp tuyến của đồ thị $y=\dfrac{x+1}{3x}$ vuông góc với $y=3x-4$ là $y=-\dfrac{1}{3}x+1$ và $y=-\dfrac{1}{3}x-\dfrac{1}{3}$. 
		\end{enumerate}
	} 
\end{ex}

% nguồn: Ok_TuanLuan_Toan11_Lan2
\begin{ex}%[1D7V1-3]%[ma:c9807f84]
	Cho hàm số $  y=x^{3}-2 x+1 $.
	\begin{enumerate}
		\item Tìm  hệ số góc của tiếp tuyến của hàm số trên tại điểm có $ x=0 $.
		\item Viết phương trình tiếp  tuyến của hàm số biết nó có $ k=-2 $.
		\item Viết phương trình tiếp tuyến của hàm số trên, biết nó tạo với  hai trục $Oxy$ một tam giác vuông cân tại $O$.
	\end{enumerate}
	 \loigiai{
		\begin{enumerate}
			 \item Ta có $ y'=3 x^{2}-2 $.\\ Hệ số góc của tiếp tuyến của hàm số tại điểm có $ x=0 $ là $ k=3\cdot (0)^2-2=-2 $.
			\item Gọi $  M\left(x_{0}; y_{0}\right) $ tiếp điểm của tiếp tuyến có hệ số góc $k=-2$.
			\[ 
			k=-2\Leftrightarrow 3  x_{0}^{2}-2=-2 \Leftrightarrow x_{0}=0.
			\] 
			 Phương trình tiếp tuyến tại điểm $ M(0; 1) $ là
			 \[
			y=-2(x-0)+1 \Leftrightarrow y=-2  x+1.
			 \]
			\item Gọi $ M\left(x_{0};y_{0}\right) $  hoành độ điểm tiếp xúc của $( C)$ và $( d)$.\\
			Gọi phương  trình đoạn chắn cắt $2$ trục tọa độ và tạo với $2$ trục $1$ tam giác vuông cân tại $O$ có dạng
			\[ 
			\dfrac{x}{a}+\dfrac{y}{b}=1  \Rightarrow y=b \cdot\left(1-\dfrac{x}{a}\right)=-\dfrac{b}{a} x+b \qquad (d)
			 \]
			Với $a \cdot b \neq 0  $; $|a|=|b|$.\\
			Khi $(d)$ là  tiếp tuyến của $ (C) $ thì $ 3 x_{0}{ }^{2}-2=\dfrac{-b}{a} $.
			Ta  có
			 \[
			|a|=|b|  \Rightarrow\hoac{&3x_0^2-2=1\\&3x_0^2-2=-1}\Rightarrow \hoac{&x_0=\pm 1\\&x_0=\pm \dfrac{\sqrt{3}}{3}.}
			\] 
			\begin{itemize}
				 \item Khi $x_0=1$ thì $y_0=0$.\\
				Phương trình tiếp tuyến là 
				\[y=1  \cdot(x-1)+0 \Leftrightarrow y=x-1.\]
				\item Khi  $x_0=-1$ thì $y_0=2$.\\
				Phương trình tiếp tuyến  là
				 \[y=1 \cdot(x+1)+2 \Leftrightarrow y=x+3.\]
				\item  Khi $x_0=\dfrac{\sqrt{3}}{3}$ thì $y_0=\dfrac{9-5\sqrt{3}}{9}$.\\
				Phương trình tiếp  tuyến là
				 \[y=-1 \cdot\left(x-\dfrac{\sqrt{3}}{3}\right)+\dfrac{9-5 \sqrt{3}}{9} \Leftrightarrow y=-x+\dfrac{9-2 \sqrt{3}}{9}.\]
				\item Khi $x_0=-\dfrac{\sqrt{3}}{3}$ thì  $y_0=\dfrac{9+5\sqrt{3}}{9}$.\\
				Phương  trình tiếp tuyến là
				\[y=-1 \cdot\left(x+\dfrac{\sqrt{3}}{3}\right)+\dfrac{9+5 \sqrt{3}}{9} \Leftrightarrow y=-x+\dfrac{9+2 \sqrt{3}}{9}.\] 
			\end{itemize}
		\end{enumerate}
	} 
\end{ex}
